\section{Một lời giải thích chỉ đúng trong một lân cận}
\label{sec:ch03-mot-loi-giai-thich-cuc-bo}

Chương trước kết thúc bằng câu hỏi về một can thiệp trên đầu vào. LIME trả lời
câu hỏi ấy bằng một \tn{mô hình thay thế cục bộ}{local surrogate}: thay vì cố
mô tả toàn bộ $f$, nó học một mô hình đơn giản $g$ cho vùng quanh một đầu vào
$x$. Ribeiro, Singh và Guestrin gọi cách làm này là local interpretable
model-agnostic explanations, hay LIME~\autocite{p09lime}.

Từ \emph{model-agnostic} có nghĩa hẹp. LIME không cần biết kiến trúc, trọng số
hay gradient của $f$; nó chỉ cần đưa một đầu vào hợp lệ vào $f$ và nhận lại đầu
ra. Do đó, cùng một khung có thể áp dụng cho bộ phân loại văn bản, bộ phân loại
ảnh hoặc một mô hình mà người giải thích không được mở bên trong. Tuy nhiên,
LIME không vì thế thoát khỏi các giả định. Nó tự chọn cách biểu diễn đầu vào,
cách tạo đầu vào đã thay đổi và cách đo khoảng cách từ mẫu ấy tới $x$.

\begin{definition}[Mô hình thay thế cục bộ]
\label{def:ch03-mo-hinh-thay-the-cuc-bo}
Cho mô hình cần giải thích $f$ và đầu vào $x$. Mô hình thay thế cục bộ là một
mô hình dễ đọc $g$ được khớp từ các đầu vào đã nhiễu quanh $x$, với mục tiêu
xấp xỉ hành vi của $f$ trong lân cận đã chọn, không phải trên mọi đầu vào.
\end{definition}

Định nghĩa~\ref{def:ch03-mo-hinh-thay-the-cuc-bo} giải thích cả sức mạnh lẫn
giới hạn của LIME. Một hàm phi tuyến phức tạp có thể trông gần như tuyến tính
trên một vùng nhỏ, nên vài trọng số của $g$ có thể giúp người đọc kiểm tra một
dự đoán. Nhưng một mô hình $g$ khớp tốt ở đầu vào này không cam kết gì về đầu
vào khác. Vì vậy, nhiều lời giải thích LIME không tự ghép thành một lời giải
thích toàn cục cho $f$.

\index{LIME|(}%
\index{mô hình thay thế cục bộ}%
